\section{Structured Failure-Aware Retrieval}
\label{sec:method}

\subsection{Pipeline}
Our repair pipeline takes buggy Python code plus a raw failure signal
(pytest output) and proceeds in five steps: (1)~extract a \emph{structured
failure state} from the failure signal; (2)~dense-retrieve candidate bug-fix
pairs from a FAISS index of corpus examples embedded with
all-MiniLM-L6-v2~\cite{sbert} over the buggy code; (3)~rerank the candidate
pool with one of the variants below; (4)~prompt the LLM with the bug, the
failure signal, and the top-$k$ ($k{=}2$) retrieved examples rendered as
compact buggy$\to$fixed diffs plus an inferred one-line ``repair lesson'';
(5)~apply the edit and run the tests. Small files are replaced whole; large
files (${>}40$k characters) are edited through a strict JSON line-range patch
schema with validation retry, which matters for repository-scale bugs
(Section~\ref{sec:results-transfer}).

\subsection{Failure state and repair metadata}
The query-side \emph{failure state} has six fields extracted from the failure
signal: \texttt{failure\_mode}, \texttt{exception\_type},
\texttt{test\_name}, \texttt{assertion\_summary},
\texttt{suspicious\_symbols}, and \texttt{contract\_tags}. On the corpus
side, each example's buggy$\to$fixed diff is mined for \emph{repair
metadata}: \texttt{repair\_pattern\_tags}, \texttt{suspicious\_symbols},
\texttt{changed\_operators}, \texttt{edit\_scope}, \texttt{file\_family}, and
\texttt{changed\_lines}.

\subsection{Retrieval variants}
We compare five variants: \textsc{baseline} (no retrieval); \codeonly (dense
retrieval on buggy code, rank = dense rank); \rawtext (dense retrieval on
buggy code + raw failure text); \rawrerank (dense retrieval + lexical
failure-text reranking); and \structured, our method, which scores each
candidate at dense rank $r$ as
\begin{align*}
\mathit{score} ={}& 1.0\cdot\mathit{dense}(r)
  + 0.8\cdot\mathit{tagOverlap}
  + 0.7\cdot\mathit{symOverlap} \\
  &+ 0.8\cdot\mathit{modeCompat}
  + 0.5\cdot\mathit{excCompat}
  + 0.2\cdot\mathit{fileFam} \\
  &- \mathit{changedLines}/400 .
\end{align*}
Overlaps are Jaccard similarities between failure-state fields and repair
metadata; compatibility terms encode hand-written failure-mode and
exception-class compatibility. The weights were fixed once, early, on
QuixBugs traces, and never tuned against any downstream result reported here.

A per-field ablation on QuixBugs (GPT-4o) confirms the signal is not
redundant: removing any one of \texttt{contract\_tags},
\texttt{suspicious\_symbols}, \texttt{failure\_mode}, \texttt{test\_name}, or
\texttt{assertion\_summary} costs one repaired problem (37/40${\to}$36/40);
only \texttt{exception\_type} contributes nothing on this benchmark.
